\chapter{NP-Hard Graph Problems and Reductions}
\label{ch:np-hard}

Through the first thirty-six chapters, every algorithm ran in polynomial time:
breadth-first search, Dijkstra, flows, matchings, and even the spectral sweeps and PageRank
iterations of the preceding chapters operated within polynomial budgets. This chapter turns
to the graph problems for which, despite decades of effort, no
polynomial-time algorithm is known, and where complexity theory suggests none exists.

What is surprising is how \emph{interconnected} these hard problems
are. Hamiltonian cycle, traveling salesman, clique, independent set, vertex cover, graph coloring,
max-cut, Steiner tree, longest path -- these look like nine unrelated puzzles. Yet in a precise
computational sense, they are polynomial-time translations of one another. A polynomial algorithm for
any one of them would give polynomial algorithms for all of them, and
for thousands of other problems besides. The equivalence comes from
\emph{reductions}: algorithms that translate one problem into another. Learning to
build and read reductions is the most useful skill for telling, in
practice, whether the problem on your desk is tractable or whether you should
stop hunting for an exact polynomial algorithm and reach for approximation,
parameterization, or heuristics.

We build the theory the way we built flows: pictures and intuition first, then
the gadgets that make reductions work, then full proofs of the main hardness
results, then the narrow boundary between easy and hard. The companion
algorithms book introduced $\mathsf{P}$, $\mathsf{NP}$, and
$\mathsf{NP}$-completeness abstractly. Here we make it concrete and graphical.

\chapterline

\section{What Hardness Means}
\label{ch37:meaning}

We work with \emph{decision problems}: questions with a yes/no answer. ``Does
$\Gph$ have a clique of size $\ge k$?'' is a decision problem; ``what is the
largest clique?'' is its optimization cousin. Decision is the right currency for
complexity theory because yes/no answers compose cleanly under reduction, and
because an optimization problem is rarely harder than polynomially many decision
queries (binary search on $k$).

\begin{definition}[The class $\mathsf{NP}$]
\label{ch37:def-np}
A decision problem $\Pi$ is in $\mathsf{NP}$ if there is a polynomial-time
\emph{verifier} $V$ and a polynomial $p$ such that for every yes-instance $x$
there exists a \emph{certificate} (witness) $w$ with $\abs{w}\le p(\abs{x})$ and
$V(x,w)=\textsf{accept}$, and for every no-instance $x$ no string $w$ makes
$V(x,w)$ accept.
\end{definition}

\begin{intuitionbox}
$\mathsf{NP}$ is the class of problems where a \emph{yes} answer has a short,
checkable proof. If I claim $\Gph$ has a Hamiltonian cycle, I hand you the cyclic
order of the vertices; you verify in linear time that consecutive vertices are
adjacent and that each vertex appears once. You do not have to find the cycle --
only check the one I show you. Finding is (apparently) hard; checking is easy.
The central question of $\mathsf{P}$ versus $\mathsf{NP}$ is whether ``easy to check''
always implies ``easy to find.''
\end{intuitionbox}

$\mathsf{P}\subseteq\mathsf{NP}$: if you can solve in polynomial time you can
verify by ignoring the certificate and re-solving. The million-dollar question is
whether the inclusion is strict. Almost everyone believes
$\mathsf{P}\neq\mathsf{NP}$, but it is unproven.

\begin{definition}[Polynomial-time many-one reduction]
\label{ch37:def-reduction}
A problem $A$ \emph{reduces} to a problem $B$, written $A\le_p B$, if there is a
polynomial-time computable function $f$ mapping instances of $A$ to instances of
$B$ such that
\[
   x \text{ is a yes-instance of } A
   \quad\Longleftrightarrow\quad
   f(x) \text{ is a yes-instance of } B.
\]
\end{definition}

Read $A\le_p B$ as ``$A$ is no harder than $B$, up to polynomial overhead.'' If
$B$ has a polynomial algorithm, compose it with $f$ to solve $A$ in polynomial
time. Contrapositive, the way we actually use it: if $A$ is hard, then $B$ is
hard too. Reductions transport hardness \emph{forward} along the arrow and
tractability \emph{backward}.

\begin{trapbox}
The direction trips up everyone at least once. To prove your new problem $B$ is
hard, you reduce a \emph{known}-hard problem $A$ \emph{to} $B$ ($A\le_p B$): you
show $B$ is at least as hard as a problem already proved $\mathsf{NP}$-hard. Reducing
$B$ to $A$ proves nothing about $B$'s hardness -- it would only show $B$ is easy
if $A$ is. Mnemonic: ``reduce \emph{from} the known hard problem \emph{to} your target.''
\end{trapbox}

\begin{definition}[$\mathsf{NP}$-hard, $\mathsf{NP}$-complete]
\label{ch37:def-npc}
A problem $B$ is \emph{$\mathsf{NP}$-hard} if $A\le_p B$ for every $A\in
\mathsf{NP}$. It is \emph{$\mathsf{NP}$-complete} if additionally $B\in
\mathsf{NP}$. The $\mathsf{NP}$-complete problems are the hardest problems in
$\mathsf{NP}$: a polynomial algorithm for any one collapses
$\mathsf{P}=\mathsf{NP}$.
\end{definition}

\begin{marginnotebox}
In 1971 Stephen Cook (and independently Leonid Levin) proved that Boolean satisfiability is $\mathsf{NP}$-complete by simulating arbitrary Turing-machine verifiers with propositional formulas. Richard Karp (1972) then proved 21 foundational combinatorial problems are $\mathsf{NP}$-complete via a web of polynomial-time reductions, establishing the standard reduction methodology.
\end{marginnotebox}

Reduction is transitive: $A\le_p B$ and $B\le_p C$ give $A\le_p C$ (compose the
two polynomial maps; a polynomial of a polynomial is a polynomial). So once
\emph{one} problem is known $\mathsf{NP}$-hard, we never again climb all the way
back to the definition. We reduce from a problem already in the club, so the club
grows one reduction at a time. We take 3-SAT's $\mathsf{NP}$-completeness as our
anchor (it is Cook-Levin plus one standard step) and grow the graph wing of the
club from there.

\begin{focusbox}
\textbf{The plan of this chapter.} Anchor at 3-SAT. Reduce
3-SAT $\to$ Independent Set $\to$ Vertex Cover $\equiv$ Clique. Separately reduce
3-SAT $\to$ 3-Coloring. Reduce 3-SAT (via a Hamiltonicity gadget) to Hamiltonian
Cycle $\to$ TSP, and to Longest Path. Touch Max-Cut and Steiner Tree. Then study
the \emph{boundary}: why 2-coloring, 2-SAT, shortest paths, and matching stay
in $\mathsf{P}$ while their natural generalizations become $\mathsf{NP}$-complete.
\end{focusbox}

\section{The Anchor: 3-SAT}
\label{ch37:3sat}

A Boolean formula in \emph{conjunctive normal form} (CNF) is an AND of clauses,
each clause an OR of literals, each literal a variable $x_i$ or its negation
$\bar x_i$. \textsc{3-SAT} asks: given a CNF formula with exactly three literals
per clause, is there an assignment of true/false to the variables making every
clause true?

\begin{theorem}[Cook-Levin, 3-SAT form]
\label{ch37:thm-3sat}
\textsc{3-SAT} is $\mathsf{NP}$-complete.
\end{theorem}

We do not reprove Cook-Levin from the Turing-machine tableau; the companion
algorithms book does that. We take it as given and note only the easy half:
\textsc{3-SAT}$\in\mathsf{NP}$ because a satisfying assignment is a certificate
checkable in linear time. From here on, every hardness proof is a \emph{gadget
construction}: turn variables and clauses into pieces of a graph.

\begin{intuitionbox}
A reduction from 3-SAT is a translation with two halves. \emph{Variable gadgets}
force a binary choice -- true or false -- one per variable. \emph{Clause gadgets}
check that at least one literal in each clause got satisfied, and wire into the
variable gadgets so they can ``see'' the choices. The key is to make the gadgets
local: each enforces its own little constraint, and the global structure of the
target problem (a large independent set, a proper coloring, a Hamiltonian cycle)
exists if and only if all the local constraints can be met simultaneously --
which is exactly satisfiability.
\end{intuitionbox}

\section{Independent Set, Vertex Cover, and Clique}
\label{ch37:isvc}

These three problems are essentially one problem. Fix a graph $\Gph=(V,E)$.

\begin{itemize}
  \item An \emph{independent set} is a set $S\subseteq V$ with no edge inside it.
        \textsc{Independent-Set}: is there one of size $\ge k$?
  \item A \emph{vertex cover} is a set $C\subseteq V$ touching every edge (each
        edge has at least one endpoint in $C$). \textsc{Vertex-Cover}: is there
        one of size $\le k$?
  \item A \emph{clique} is a set $K\subseteq V$ that is pairwise adjacent.
        \textsc{Clique}: is there one of size $\ge k$?
\end{itemize}

\begin{lemma}[The complementarity triangle]
\label{ch37:lem-triangle}
For any graph $\Gph$ on $n$ vertices and any $S\subseteq V$:
\begin{enumerate}
  \item $S$ is independent in $\Gph$ $\iff$ $V\setminus S$ is a vertex cover of
        $\Gph$.
  \item $S$ is independent in $\Gph$ $\iff$ $S$ is a clique in the complement
        $\bar\Gph$.
\end{enumerate}
Hence $\Gph$ has an independent set of size $\ge k$ iff it has a vertex cover of
size $\le n-k$ iff $\bar\Gph$ has a clique of size $\ge k$.
\end{lemma}

\begin{proof}
(1) $S$ independent means no edge has both endpoints in $S$, i.e. every edge has
an endpoint outside $S$, i.e. $V\setminus S$ covers every edge -- and the chain
of ``i.e.''s is reversible. (2) An edge inside $S$ in $\Gph$ is a non-edge inside
$S$ in $\bar\Gph$, so $S$ edge-free in $\Gph$ $\iff$ $S$ complete in $\bar\Gph$.
The size statements follow since $\abs{V\setminus S}=n-\abs S$, and complementing
the graph is a polynomial-time, instance-preserving map.
\end{proof}

So these three are interreducible by trivial transformations: complement the
graph, or complement the set. Figure~\ref{ch37:fig-triangle-duality} shows both
transformations side by side on a four-vertex paw graph. Prove one hard and all
three are hard. We prove \textsc{Independent-Set} hard by reducing from 3-SAT.

\begin{figure}[H]
\centering
\begin{tikzpicture}[
  vx/.style={circle,draw=edgegray,fill=white,thick,minimum size=16pt,inner sep=1pt,font=\scriptsize\bfseries},
  isvx/.style={circle,draw=sourceblue,fill=sourceblue!25,very thick,minimum size=16pt,inner sep=1pt,font=\scriptsize\bfseries},
  vcvx/.style={circle,draw=quotientorange,fill=quotientorange!25,very thick,minimum size=16pt,inner sep=1pt,font=\scriptsize\bfseries},
  clvx/.style={circle,draw=targetgreen,fill=targetgreen!25,very thick,minimum size=16pt,inner sep=1pt,font=\scriptsize\bfseries}]

  % Panel 1: Independent Set S in G
  \begin{scope}[shift={(0,0)}]
    \draw[rounded corners=6pt,fill=sourceblue!4,draw=sourceblue!25,thick] (-1.2,-1.15) rectangle (2.3,1.15);
    \node[font=\footnotesize\bfseries,text=sourceblue] at (0.55,1.45) {1.\ Independent Set $S$ in $\Gph$};
    \node[isvx] (a1) at (-0.5,0.55) {$1$};
    \node[vx]   (a2) at (-0.5,-0.55) {$2$};
    \node[vx]   (a3) at (0.6,0) {$3$};
    \node[isvx] (a4) at (1.7,0) {$4$};
    \draw[edgegray,thick] (a1)--(a2)--(a3)--(a1) (a3)--(a4);
    \node[font=\scriptsize,text=sourceblue] at (0.55,-1.45) {$S=\{1,4\}$ independent ($\abs{S}=2$)};
  \end{scope}

  % Panel 2: Vertex Cover V \ S in G
  \begin{scope}[shift={(4.4,0)}]
    \draw[rounded corners=6pt,fill=quotientorange!4,draw=quotientorange!30,thick] (-1.2,-1.15) rectangle (2.3,1.15);
    \node[font=\footnotesize\bfseries,text=quotientorange] at (0.55,1.45) {2.\ Vertex Cover $V\setminus S$ in $\Gph$};
    \node[vx]   (b1) at (-0.5,0.55) {$1$};
    \node[vcvx] (b2) at (-0.5,-0.55) {$2$};
    \node[vcvx] (b3) at (0.6,0) {$3$};
    \node[vx]   (b4) at (1.7,0) {$4$};
    \draw[quotientorange,thick] (b1)--(b2)--(b3)--(b1) (b3)--(b4);
    \node[font=\scriptsize,text=quotientorange] at (0.55,-1.45) {$V\setminus S=\{2,3\}$ covers all edges};
  \end{scope}

  % Panel 3: Clique S in complement G-bar
  \begin{scope}[shift={(8.8,0)}]
    \draw[rounded corners=6pt,fill=targetgreen!4,draw=targetgreen!30,thick] (-1.2,-1.15) rectangle (2.3,1.15);
    \node[font=\footnotesize\bfseries,text=targetgreen] at (0.55,1.45) {3.\ Clique $S$ in $\overline{\Gph}$};
    \node[clvx] (c1) at (-0.5,0.55) {$1$};
    \node[vx]   (c2) at (-0.5,-0.55) {$2$};
    \node[vx]   (c3) at (0.6,0) {$3$};
    \node[clvx] (c4) at (1.7,0) {$4$};
    \draw[targetgreen,very thick] (c1) to[bend left=22] (c4);
    \draw[edgegray,thick,dashed] (c2) to[bend right=22] (c4);
    \node[font=\scriptsize,text=targetgreen] at (0.55,-1.45) {$S=\{1,4\}$ adjacent in $\overline{\Gph}$ ($\abs{K}=2$)};
  \end{scope}
\end{tikzpicture}
\caption{The complementarity triangle (Lemma~\ref{ch37:lem-triangle}) on a $4$-vertex graph $\Gph$. Taking the vertex complement turns the independent set $S=\{1,4\}$ into the vertex cover $V\setminus S=\{2,3\}$ of $\Gph$; taking the edge complement turns the same set $S=\{1,4\}$ into a clique in $\overline{\Gph}$.}
\label{ch37:fig-triangle-duality}
\end{figure}

\subsection{The gadget}

\begin{theorem}
\label{ch37:thm-is}
\textsc{Independent-Set} is $\mathsf{NP}$-complete.
\end{theorem}

\begin{proof}
Membership: a set of $k$ pairwise-nonadjacent vertices is a certificate checkable
in $\Oh(k^2)$ time. Hardness: reduce \textsc{3-SAT}. Let $\varphi$ have $m$
clauses $C_1,\dots,C_m$, each with three literals. Build a graph $\Gph_\varphi$:

\begin{itemize}
  \item \textbf{Clause triangles.} For each clause $C_j=(\ell_{j1}\vee
        \ell_{j2}\vee \ell_{j3})$ create three vertices, one per literal
        occurrence, and join them in a triangle. A triangle admits at most one
        vertex in any independent set -- this forces us to pick exactly one
        ``satisfying literal'' per clause.
  \item \textbf{Conflict edges.} For every pair of literal-vertices that are
        \emph{negations} of each other ($x_i$ in one place, $\bar x_i$ in
        another), add an edge. This forbids choosing both $x_i$ and $\bar x_i$ --
        it keeps our choices consistent with a single assignment.
\end{itemize}

Set $k=m$. We claim $\varphi$ is satisfiable $\iff$ $\Gph_\varphi$ has an
independent set of size $m$.
\end{proof}

\begin{figure}[H]
\centering
\begin{tikzpicture}[
  lit/.style={circle,draw=nodeblue,fill=nodeblue!12,thick,inner sep=1pt,font=\scriptsize\bfseries,minimum size=17pt},
  chosen/.style={circle,draw=sourceblue,fill=sourceblue!30,very thick,inner sep=1pt,font=\scriptsize\bfseries,minimum size=17pt},
  clause/.style={draw=edgegray,thick},
  conflict/.style={draw=trapred,very thick,dashed}]

  % Clause 1 panel
  \draw[rounded corners=8pt,fill=sourceblue!5,draw=sourceblue!25,thick] (-1.2,-0.2) rectangle (1.2,2.7);
  \node[font=\footnotesize\bfseries,text=sourceblue] at (0,2.95) {Clause $C_1 = (x_1 \vee x_2 \vee x_3)$};

  \node[chosen] (a1) at (0,2.1)     {$x_1$};
  \node[lit]    (a2) at (-0.7,0.6)  {$x_2$};
  \node[lit]    (a3) at (0.7,0.6)   {$x_3$};
  \draw[clause] (a1)--(a2)--(a3)--(a1);

  % Clause 2 panel
  \draw[rounded corners=8pt,fill=targetgreen!5,draw=targetgreen!25,thick] (3.8,-0.2) rectangle (6.2,2.7);
  \node[font=\footnotesize\bfseries,text=targetgreen] at (5.0,2.95) {Clause $C_2 = (\bar x_1 \vee \bar x_2 \vee x_3)$};

  \node[lit]    (b1) at (5.0,2.1)   {$\bar x_1$};
  \node[lit]    (b2) at (4.3,0.6)   {$\bar x_2$};
  \node[chosen] (b3) at (5.7,0.6)   {$x_3$};
  \draw[clause] (b1)--(b2)--(b3)--(b1);

  % Conflict edges between complementary literals
  \draw[conflict] (a1) to[bend left=16] node[above=2pt,font=\tiny\bfseries,text=trapred]{conflict} (b1);
  \draw[conflict] (a2) to[bend left=12] node[below=2pt,font=\tiny\bfseries,text=trapred]{conflict} (b2);

  % Legend / note
  \node[font=\scriptsize\itshape,text=edgegray,align=center] at (2.5,-0.6)
    {Solid triangles enforce at most 1 chosen literal per clause.\\Dashed red edges forbid contradictory assignments ($x_i \wedge \bar{x}_i$).};
\end{tikzpicture}
\caption{The 3-SAT to Independent Set reduction for $(x_1\vee x_2\vee x_3)\wedge(\bar x_1\vee \bar x_2\vee x_3)$.
Each clause is a solid triangle requiring at least two vertices to be covered (hence at most one in the independent set).
Dashed conflict edges connect complementary literals. An independent set of size $m=2$ (shaded blue: $x_1$ and $x_3$)
corresponds to a consistent satisfying assignment (setting $x_1=\text{true}, x_3=\text{true}$).}
\label{ch37:fig-is}
\end{figure}

\begin{whybox}
\textbf{Why the construction is correct.}
($\Rightarrow$) Take a satisfying assignment. Each clause has at least one true
literal; pick one such literal-vertex per triangle. That is $m$ vertices, one per
triangle, so no triangle edge is violated. Two chosen vertices can only conflict
via a dashed edge, which connects $x_i$ to $\bar x_i$ -- but both being chosen
would mean both literals are true under the same assignment, impossible. So the
$m$ vertices are independent.

($\Leftarrow$) Take an independent set $S$ of size $m$. Each triangle holds at
most one vertex of $S$, and there are $m$ triangles, so $S$ holds \emph{exactly}
one vertex per triangle. Define an assignment by setting each chosen literal true
(if $x_i$'s vertex is chosen set $x_i$ true; if $\bar x_i$'s vertex is chosen set
$x_i$ false). No conflict edge is internal to $S$, so we never demand $x_i$ both
true and false -- the assignment is well defined. Variables not pinned by any
chosen literal get an arbitrary value. Every clause has its chosen literal true,
so $\varphi$ is satisfied.
\end{whybox}

The reduction builds $3m$ vertices and at most $3m + \binom{3m}{2}$ edges in
$\Oh(m^2)$ time -- polynomial. By Lemma~\ref{ch37:lem-triangle},
\textsc{Vertex-Cover} and \textsc{Clique} are $\mathsf{NP}$-complete too.

\begin{labbox}
Trace the formula $\varphi=(x_1\vee x_2\vee x_3)\wedge(\bar x_1\vee\bar
x_2\vee x_3)$ from Figure~\ref{ch37:fig-is}. We want an independent set of size
$2$. Try the assignment $x_3=\text{true}$ (with $x_1,x_2$ free). Pick the $x_3$
vertex in $C_1$ and the $x_3$ vertex in $C_2$: no triangle edge between them, no
conflict edge (both are the literal $x_3$, not negations), so $\{x_3^{C_1},
x_3^{C_2}\}$ is independent of size $2$. The reduction confirms $\varphi$ is
satisfiable (indeed $x_3=$true satisfies both clauses). Now try to force a
conflict: pick $x_1$ in $C_1$ and $\bar x_1$ in $C_2$ -- the dashed edge between
them blocks it, exactly as a consistent assignment must.
\end{labbox}

\begin{algorithm}[H]
\caption{\textsc{Reduce-3SAT-to-IndependentSet}$(\varphi)$}
\label{ch37:alg-is}
\DontPrintSemicolon
\KwIn{CNF formula $\varphi$ with clauses $C_1,\dots,C_m$, three literals each.}
\KwOut{Graph $\Gph$ and integer $k$ with $\varphi$ sat $\iff$ $\Gph$ has IS of size $k$.}
$V\gets\emptyset$;\quad $E\gets\emptyset$\;
\ForEach{clause $C_j=(\ell_{j1}\vee\ell_{j2}\vee\ell_{j3})$}{
  create vertices $v_{j1},v_{j2},v_{j3}$ labeled by the literals\;
  add triangle edges $\{v_{j1}v_{j2}, v_{j2}v_{j3}, v_{j3}v_{j1}\}$ to $E$\;
}
\ForEach{pair $u,v$ with $\mathrm{label}(u)=x_i,\ \mathrm{label}(v)=\bar x_i$}{
  add edge $uv$ to $E$ \tcp*{conflict edge}
}
$k\gets m$\;
\Return $(\Gph=(V,E),\,k)$\;
\end{algorithm}

\begin{complexitybox}
The map runs in $\Oh(m^2)$ time and produces $\Oh(m)$ vertices, $\Oh(m^2)$ edges.
What it buys: \textsc{Independent-Set}, \textsc{Vertex-Cover}, and
\textsc{Clique} are $\mathsf{NP}$-complete. No polynomial exact algorithm is
known; the textbook brute force is $\Oh(2^n\cdot\poly(n))$, improvable to roughly
$\Oh(1.2^n)$ for \textsc{Independent-Set} by clever branching.
\end{complexitybox}

\section{Graph Coloring}
\label{ch37:coloring}

A \emph{proper $c$-coloring} assigns each vertex one of $c$ colors so that
adjacent vertices differ. \textsc{$c$-Coloring}: can $\Gph$ be properly colored
with $c$ colors? The chromatic number $\chi(\Gph)$ is the least feasible $c$.

The jump from $c=2$ to $c=3$ is the cleanest easy/hard boundary in the subject,
so it earns its own section (\S\ref{ch37:boundary}). Here we prove the hard side.

\begin{theorem}
\label{ch37:thm-3col}
\textsc{3-Coloring} is $\mathsf{NP}$-complete.
\end{theorem}

\begin{proof}[Construction]
Membership: a coloring is a certificate, checked by scanning edges. Hardness:
reduce \textsc{3-SAT}. Colors are named $\mathsf{T}$, $\mathsf{F}$, $\mathsf{B}$
(``base''). The reduction uses three kinds of pieces.

\textbf{Palette triangle.} Three special vertices $T,F,B$ joined in a triangle.
In any 3-coloring they receive all three colors; rename so $T$ gets $\mathsf{T}$,
$F$ gets $\mathsf{F}$, $B$ gets $\mathsf{B}$. They anchor the meaning of the
colors.

\textbf{Variable gadget.} For each variable $x_i$, two vertices $x_i$ and $\bar
x_i$ joined to each other and both joined to $B$. The triangle $x_i,\bar x_i,B$
forces $x_i$ and $\bar x_i$ to take the two colors $\{\mathsf{T},\mathsf{F}\}$ in
some order -- exactly a Boolean choice, with $x_i$ and $\bar x_i$ guaranteed
opposite.

\textbf{Clause gadget (OR gadget).} For a clause $(a\vee b\vee c)$ attach a small
widget that is 3-colorable \emph{iff} at least one of $a,b,c$ is colored
$\mathsf{T}$. The widget is the classic ``OR tree'': two binary OR-gadgets
chained, the output forced adjacent to $F$ and $B$. A binary OR-gadget on inputs
$p,q$ adds three vertices: $u_p$ adjacent to $p$, $u_q$ adjacent to $q$, and an
output $r$, with $u_p,u_q,r$ forming a triangle.
\end{proof}

\begin{figure}[H]
\centering
\begin{tikzpicture}[every node/.style={circle,draw,thick,inner sep=1.5pt,
  font=\scriptsize,minimum size=15pt},scale=0.95]
\node[fill=green!18] (T) at (-3.8,3.1) {$T$};
\node[fill=red!18]   (F) at (-3.8,1.5) {$F$};
\node[fill=gray!25]  (B) at (-2.7,2.3) {$B$};
\draw[edgegray,thick] (T)--(F)--(B)--(T);
\node (xi)  at (-0.6,3.0) {$x_i$};
\node (nxi) at (-0.6,1.6) {$\bar x_i$};
\draw[edgegray,thick] (xi)--(nxi) (B)--(xi) (B)--(nxi);
\node (a)   at (1.8,3.3) {$a$};
\node (b)   at (1.8,1.9) {$b$};
\node (o1)  at (3.0,2.6) {};
\node (o2)  at (4.1,2.6) {};
\node (c)   at (3.5,3.6) {$c$};
\node (out) at (5.2,2.6) {};
\draw[edgegray,thick] (a)--(o1) (b)--(o1) (o1)--(o2) (c)--(o2) (o2)--(out);
\draw[edgegray,thick] (out) to[out=-125,in=-45] (F);
\draw[edgegray,thick] (out) to[out=125,in=55] (B);
\node[draw=none,font=\scriptsize] at (-3.3,-0.3) {palette};
\node[draw=none,font=\scriptsize] at (-0.6,-0.3) {variable $x_i$};
\node[draw=none,font=\scriptsize] at (3.6,-0.3) {clause $(a\vee b\vee c)$ OR-gadget};
\end{tikzpicture}
\caption{3-coloring gadgets. The palette triangle fixes the three colors. Each
variable's $x_i,\bar x_i$ are forced opposite by the triangle with $B$. The
clause part is schematic: each unlabeled node stands for the output of one
binary OR-gadget, whose internal triangle is not drawn. The
OR-gadget output can be colored $\mathsf{T}$ only when some input is
$\mathsf{T}$; tying the output to $F$ and $B$ forbids the all-$\mathsf{F}$ case.}
\label{ch37:fig-3col}
\end{figure}

\begin{whybox}
\textbf{Why the OR-gadget works.} The atomic OR-gadget is a small graph on inputs
$p,q$ and an output node $r$ with the property: in any 3-coloring with palette
fixed, $r$ can be colored $\mathsf{T}$ iff $p$ or $q$ is $\mathsf{T}$, and if both
$p,q$ are $\mathsf{F}$ then $r$ is forced to $\mathsf{F}$. Check it on the triangle
$u_p,u_q,r$. If $p=q=\mathsf{F}$, then $u_p$ and $u_q$ avoid $\mathsf{F}$ and each
other, so they take $\mathsf{T}$ and $\mathsf{B}$, and $r$ is left with
$\mathsf{F}$. If $p=\mathsf{T}$, color $u_p=\mathsf{F}$ and $u_q=\mathsf{B}$ (legal
whether $q$ is $\mathsf{T}$ or $\mathsf{F}$), and $r=\mathsf{T}$ is free; the case
$q=\mathsf{T}$ is symmetric. Chaining two of them
computes the OR of three literals. Finally we make the clause's output vertex
adjacent to both $F$ (color $\mathsf{F}$) and $B$ (color $\mathsf{B}$), which
forbids the output from being $\mathsf{F}$ or $\mathsf{B}$ -- so it must be
$\mathsf{T}$, which is possible exactly when some literal of the clause is
$\mathsf{T}$, i.e. true. A formula is satisfiable iff every clause-output can be
forced $\mathsf{T}$ iff the whole graph is 3-colorable.
\end{whybox}

The construction has $\Oh(n+m)$ vertices and edges and is built in linear time.
\textsc{3-Coloring} is $\mathsf{NP}$-complete, and for every fixed $c\ge 3$,
\textsc{$c$-Coloring} is $\mathsf{NP}$-complete (add a $(c-3)$-clique joined to
the whole graph to lift 3 to $c$).

\begin{applicationbox}
Register allocation in compilers is graph coloring: vertices are program
variables alive at the same time, edges mean ``cannot share a register,'' colors
are physical registers. Frequency assignment for radio towers, exam timetabling
(colors are time slots, edges are conflicts), and Sudoku (extending a pre-colored
81-vertex conflict graph to a proper 9-coloring) are all direct instances of graph coloring. Because exact coloring is
$\mathsf{NP}$-hard, real compilers use the greedy + spilling heuristics of
Chaitin-Briggs rather than optimal colorings.
\end{applicationbox}

\section{Hamiltonian Cycle and the Traveling Salesman}
\label{ch37:ham-tsp}

A \emph{Hamiltonian cycle} visits every vertex exactly once and returns to start.
\textsc{Ham-Cycle}: does $\Gph$ have one? Compare it with the
\emph{Eulerian} cycle (visit every \emph{edge} once), which is decidable in
linear time -- a graph is Eulerian iff it is connected with all degrees even.
Swapping ``edge'' for ``vertex'' turns a linear-time problem into an
$\mathsf{NP}$-complete one. The local-to-global structure that made Euler easy
(degree parity is a purely local certificate) has no analogue for Hamiltonicity.

\begin{marginnotebox}
The Hamiltonian cycle is named for William Rowan Hamilton's 1857 ``Icosian Game'' on a dodecahedron. Euler's 1736 bridge problem (visiting every edge) is solvable in linear time, but Hamilton's problem (visiting every vertex) is $\mathsf{NP}$-complete.
\end{marginnotebox}

\begin{theorem}
\label{ch37:thm-ham}
\textsc{Ham-Cycle} (directed and undirected) is $\mathsf{NP}$-complete.
\end{theorem}

\begin{proof}[Sketch with the key gadget]
Membership: the cyclic vertex order is a linear-time-checkable certificate.
Hardness: reduce \textsc{3-SAT} (directed version; an undirected version follows
by replacing each directed vertex with a 3-vertex path that enforces direction).
The reduction uses a \emph{choice-and-check} layout.

\textbf{Variable chains.} For each variable $x_i$ build a horizontal chain of $2m$
vertices that the Hamiltonian cycle must traverse. Traversing it left-to-right
encodes $x_i=\text{true}$; right-to-left encodes $x_i=\text{false}$. The chains
are stacked and linked into one big loop, so the cycle threads each chain in one
of two directions -- a global binary choice per variable.

\textbf{Clause vertices.} For each clause $C_j$ add a single vertex $c_j$. Wire
$c_j$ to a position on the chain of each variable appearing in $C_j$, oriented so
that the cycle can detour into $c_j$ and back \emph{only} when the corresponding
chain is traversed in the direction that satisfies the literal. Thus $c_j$ can be
visited iff at least one of its literals is true.
\end{proof}

\begin{figure}[H]
\centering
\begin{tikzpicture}[>=Stealth,
  vNode/.style={circle,draw=nodeblue,fill=nodeblue!15,thick,minimum size=6.5mm,font=\tiny\bfseries,inner sep=0pt},
  cNode/.style={circle,draw=quotientorange,fill=quotientorange!20,thick,minimum size=7.5mm,font=\scriptsize\bfseries,inner sep=0pt},
  chainArr/.style={->,targetgreen,very thick},
  detourArr/.style={->,trapred,very thick}]

  % Panel frame
  \draw[rounded corners=6pt,fill=sourceblue!4,draw=sourceblue!25,thick] (-1.2,-2.2) rectangle (8.6,2.2);

  % Variable chain x_i (top lane)
  \node[font=\footnotesize\bfseries,text=sourceblue,anchor=west] at (-0.9,1.8) {Variable Chain $x_i$: Left-to-Right = \textsc{True}, Right-to-Left = \textsc{False}};

  \node[vNode] (x1) at (0,0.8) {$v_1$};
  \node[vNode] (x2) at (1.5,0.8) {$v_2$};
  \node[vNode] (x3) at (3.0,0.8) {$v_3$};
  \node[vNode] (x4) at (4.5,0.8) {$v_4$};
  \node[vNode] (x5) at (6.0,0.8) {$v_5$};
  \node[vNode] (x6) at (7.5,0.8) {$v_6$};

  \draw[chainArr] (x1) -- (x2);
  \draw[chainArr] (x2) -- (x3);
  \draw[chainArr] (x4) -- (x5);
  \draw[chainArr] (x5) -- (x6);

  % Detour to clause vertex c_j
  \node[cNode] (cj) at (3.75,-0.6) {$c_j$};
  \node[font=\scriptsize\bfseries,text=quotientorange,below=2pt] at (cj.south) {Clause $C_j$};

  \draw[detourArr] (x3) to[bend right=20] node[left=2pt,font=\tiny\bfseries,text=trapred]{detour in} (cj);
  \draw[detourArr] (cj) to[bend right=20] node[right=2pt,font=\tiny\bfseries,text=trapred]{detour out} (x4);

  % Direct edge between x3 and x4
  \draw[->,edgegray,dashed,thick] (x3) -- node[above=1pt,font=\tiny]{direct edge} (x4);

  % Explanation text at bottom
  \node[font=\scriptsize\itshape,text=edgegray,align=center] at (3.75,-1.7)
    {When $x_i=\text{true}$, the cycle enters $v_3 \to c_j \to v_4$, visiting $c_j$ on the detour.\\If traversed in reverse ($x_i=\text{false}$), the directed detour cannot be taken.};
\end{tikzpicture}
\caption{The Hamiltonian path/cycle gadget for 3-SAT. Each variable $x_i$ is mapped to a horizontal
directed chain. If $x_i$ is set to true, the cycle moves left-to-right and can take the detour $v_3 \to c_j \to v_4$
to visit the clause vertex $c_j$. The figure draws each chain step in one direction only; the full gadget has
arcs both ways between consecutive chain vertices (with spacer vertices between clause positions), so the cycle
can run along the chain in either direction. If the chain is traversed right-to-left ($x_i=\text{false}$), the
one-way clause arcs prevent taking this detour. Thus, $c_j$ can be visited without breaking the cycle if and only if at least
one of its literals evaluates to true.}
\label{ch37:fig-ham}
\end{figure}

Hamiltonicity leads to the best-known hard optimization problem.

\begin{definition}[Traveling Salesman]
\label{ch37:def-tsp}
\textsc{TSP}: given $n$ cities and pairwise distances $d(i,j)$ and a bound $L$, is
there a tour visiting every city once with total length $\le L$?
\end{definition}

\begin{theorem}
\label{ch37:thm-tsp}
\textsc{TSP} is $\mathsf{NP}$-complete, even in the \emph{metric} case where
distances obey the triangle inequality.
\end{theorem}

\begin{proof}
Reduce \textsc{Ham-Cycle}. Given $\Gph=(V,E)$, build a complete graph on $V$ with
\[
  d(i,j)=\begin{cases} 1 & \{i,j\}\in E,\\ 2 & \{i,j\}\notin E,\end{cases}
\]
and set the bound $L=n$. Distances are $1$ or $2$, so the triangle inequality
holds automatically ($d(i,k)\le 2\le d(i,j)+d(j,k)$), giving the metric case for
free. A tour visits $n$ edges. Its length is $n$ iff every edge it uses has
weight $1$ iff every consecutive pair is a real edge of $\Gph$ iff the tour is a
Hamiltonian cycle of $\Gph$. So $\Gph$ is Hamiltonian $\iff$ the TSP instance has
a tour of length $\le n$. The map is $\Oh(n^2)$.
\end{proof}

\begin{trapbox}
\textsc{Ham-Cycle} reduces to \textsc{TSP} the easy way (weights $1$/$2$), which
\emph{also} settles metric TSP. But \emph{general} TSP -- arbitrary weights, no
triangle inequality -- is much worse to approximate: encode Hamiltonicity with
weight $1$ for edges and weight $n\cdot\rho$ for non-edges, and a tour shorter
than $n\rho$ exists iff $\Gph$ is Hamiltonian. So a polynomial
$\rho$-approximation to general TSP would decide Hamiltonicity -- general TSP has
\emph{no} constant-factor approximation unless $\mathsf{P}=\mathsf{NP}$. Metric
TSP, by contrast, has the Christofides $3/2$-approximation. The triangle
inequality is what separates them; do not conflate the two problems.
\end{trapbox}

\section{Longest Simple Path}
\label{ch37:longest}

Shortest path was solved in polynomial time in Chapters 7--9: BFS, Dijkstra, Bellman-Ford,
Floyd-Warshall, all polynomial. \textsc{Longest-Path} -- given $\Gph$ and $k$, is
there a \emph{simple} path with $\ge k$ edges? -- is $\mathsf{NP}$-complete.

\begin{theorem}
\label{ch37:thm-longest}
\textsc{Longest-Path} is $\mathsf{NP}$-complete, even on unweighted graphs.
\end{theorem}

\begin{proof}
Membership: a path is a checkable certificate. Hardness: reduce
\textsc{Ham-Path} (a Hamiltonian path, the open-ended cousin of Ham-Cycle, is
also $\mathsf{NP}$-complete by the same gadget without closing the loop). A graph
on $n$ vertices has a Hamiltonian path iff it has a simple path with $\ge n-1$
edges. Set $k=n-1$; the identity map is the reduction.
\end{proof}

\begin{whybox}
\textbf{Why does ``longest'' break what ``shortest'' kept?} The difference is the word
\emph{simple}. Shortest paths are automatically simple -- a cycle only adds
nonnegative length, so optimal solutions never repeat a vertex, and that lets
dynamic programming over distances work: the prefix of a shortest path is a
shortest path (optimal substructure). Longest \emph{simple} paths have no such
substructure: the longest path from $s$ to $v$ may not extend the longest path
from $s$ to $u$, because reusing vertices is forbidden and the optimal prefix may
``waste'' a vertex you needed later. The moment you drop simplicity to recover
substructure, the longest \emph{walk} blows up to $+\infty$ whenever a cycle is
reachable (every traversal of it adds length), so the walk relaxation answers a
different, useless question. Either way DP fails. Negative cycles break shortest paths for exactly
the mirror-image reason; the two pathologies are dual.
\end{whybox}

\begin{intuitionbox}
A simple way to see the boundary: longest path on a DAG \emph{is} polynomial
(topological-order DP), because a DAG has no cycles to repeat -- simplicity is
free. So \textsc{Longest-Path} is easy on DAGs and hard on general graphs. The
hardness comes from the interaction between cycles and the simplicity
constraint.
\end{intuitionbox}

\section{Max-Cut and Steiner Tree}
\label{ch37:maxcut-steiner}

\textbf{Max-Cut.} Partition $V$ into $(S,\bar S)$ to maximize the number of edges
crossing. Decision: is there a cut of size $\ge k$? Its minimization counterpart,
\emph{Min-Cut} (Chapters 11--12), is polynomial via max-flow.

\begin{theorem}
\label{ch37:thm-maxcut}
\textsc{Max-Cut} is $\mathsf{NP}$-complete.
\end{theorem}

\begin{proof}[Idea]
Reduce from \textsc{Not-All-Equal-3-SAT} (NAE-3SAT), a known $\mathsf{NP}$-complete variant of 3-SAT.
In NAE-3SAT, a clause is satisfied if its three literals do not all take the same truth value -- that is,
they cannot be all-true or all-false. This requirement translates directly into a cut: think of the two
sides of the cut $(S, \bar{S})$ as representing truth values $\text{true}$ and $\text{false}$.
Make one vertex per literal, shared by all clauses: $x_i$ and $\bar x_i$ for each of the $n$ variables.
Each clause adds a triangle on its three literal vertices (parallel edges allowed). A clause triangle cuts at
most $2$ edges, and it cuts \emph{exactly} $2$ edges if and only if the three literals are not all placed on
the same side (Figure~\ref{ch37:fig-maxcut-gadget}). To make the two sides a consistent assignment, join each
pair $x_i,\bar x_i$ by $M=2m+1$ parallel edges, where $m$ is the number of clauses. A cut that separates
every pair earns $nM$ from these edges; a cut that fails to separate even one pair has at most
$(n-1)M+2m<nM$ in total. So a maximum cut separates every pair, and its value reaches $nM+2m$ exactly when
every clause triangle contributes $2$, which happens precisely when a valid not-all-equal truth assignment
exists.
\end{proof}

\begin{figure}[H]
\centering
\begin{tikzpicture}[scale=1.1,
  vTrue/.style={circle,draw=sourceblue,fill=sourceblue!15,thick,minimum size=7mm,font=\scriptsize\bfseries},
  vFalse/.style={circle,draw=targetgreen,fill=targetgreen!15,thick,minimum size=7mm,font=\scriptsize\bfseries},
  triEdge/.style={draw=edgegray,thick},
  cutLine/.style={draw=trapred,very thick,dashed}]

  % ================= Left Panel: All-Equal =================
  \begin{scope}[shift={(0,0)}]
    \draw[rounded corners=6pt,fill=gray!4,draw=gray!30,thick] (-1.4,-0.6) rectangle (1.4,2.8);
    \node[font=\footnotesize\bfseries,text=trapred] at (0,2.5) {All Equal (Unsatisfied NAE)};

    \node[vTrue] (a1) at (0,1.5) {$\ell_1$};
    \node[vTrue] (a2) at (-0.8,0.2) {$\ell_2$};
    \node[vTrue] (a3) at (0.8,0.2) {$\ell_3$};

    \draw[triEdge] (a1) -- (a2) -- (a3) -- (a1);

    \node[font=\scriptsize\bfseries,text=trapred,below=4pt] at (0,-0.6) {Edges Cut $= 0$};
  \end{scope}

  % ================= Right Panel: Not-All-Equal =================
  \begin{scope}[shift={(5.2,0)}]
    \draw[rounded corners=6pt,fill=sourceblue!4,draw=sourceblue!25,thick] (-1.8,-0.6) rectangle (1.8,2.8);
    \node[font=\footnotesize\bfseries,text=targetgreen!80!black] at (0,2.5) {Not-All-Equal (Satisfied NAE)};

    \node[vTrue] (b1) at (-0.7,1.5) {$\ell_1$};
    \node[vTrue] (b2) at (-0.7,0.2) {$\ell_2$};
    \node[vFalse] (b3) at (0.8,0.85) {$\ell_3$};

    \draw[triEdge] (b1) -- (b2);
    \draw[triEdge] (b1) -- (b3);
    \draw[triEdge] (b2) -- (b3);

    % Cut dividing left (True) and right (False)
    \draw[cutLine] (0.05,-0.5) -- (0.05,1.9)
      node[above=2pt,font=\tiny\bfseries,text=trapred]{Cut $(S \mid \bar{S})$};

    \node[font=\scriptsize\bfseries,text=targetgreen!80!black,below=4pt] at (0,-0.6) {Edges Cut $= 2$ (Maximum possible)};
  \end{scope}

\end{tikzpicture}
\caption{The clause triangle gadget in the reduction from NAE-3SAT to Max-Cut. Each clause is represented
by a triangle of literal vertices. Left: If all three literals take the same truth value, they lie on the
same side of the partition and $0$ edges are cut. Right: If the clause is not-all-equal, the cut separates
one literal from the other two, cutting exactly $2$ edges. Maximizing the total cut forces every clause
to achieve this maximum.}
\label{ch37:fig-maxcut-gadget}
\end{figure}

\begin{connectionbox}
Max-Cut connects to two later chapters. Its \emph{spectral} relaxation --
maximize $\tfrac14 x^\top \Lap x$ over $x\in\{\pm1\}^n$, where $\Lap$ is the graph
Laplacian (Chapter 31) -- connects cut size to the eigenvalues of $\Lap$ and to
Cheeger-type bounds (Chapter 32). Relaxing $\{\pm1\}^n$ to unit vectors gives the
Goemans-Williamson semidefinite program and its famous $0.878$-approximation
(Chapter 38). So the same problem that is $\mathsf{NP}$-hard here becomes a question
about Laplacian quadratic forms and rounding hyperplanes there.
\end{connectionbox}

\textbf{Steiner Tree.} Given a graph with edge weights and a subset $R\subseteq V$
of \emph{terminals}, find a minimum-weight connected subgraph containing all of
$R$. Decision: weight $\le L$? When $R=V$ this is the (polynomial) minimum
spanning tree; when $\abs R=2$ it is shortest path. The general case -- choose
which non-terminal ``Steiner'' vertices to include -- is $\mathsf{NP}$-complete,
by reduction from \textsc{Exact-Cover-by-3-Sets} or from \textsc{Vertex-Cover}.

\begin{intuitionbox}
Steiner tree shows that the hardness is not about ``all'' or ``two'' but about the
combinatorial choice in between: which optional vertices to recruit. MST (recruit
everyone) and shortest path (recruit no one) are the easy endpoints; the
hardness comes from the exponentially many intermediate recruitment choices
(Figure~\ref{ch37:fig-steiner}).
There is a metric $\ln 4 + \eps \approx 1.39$-approximation (Byrka et al.) and a
classic $2$-approximation via the terminal metric's MST.
\end{intuitionbox}

\begin{figure}[H]
\centering
\begin{tikzpicture}[every node/.style={circle,draw,thick,inner sep=1.5pt,
  font=\scriptsize,minimum size=15pt},scale=1.0]
\node[fill=sourceblue!25] (t1) at (-1.4,1.4) {$t_1$};
\node[fill=sourceblue!25] (t2) at (1.4,1.4) {$t_2$};
\node[fill=sourceblue!25] (t3) at (1.4,-1.4) {$t_3$};
\node[fill=sourceblue!25] (t4) at (-1.4,-1.4) {$t_4$};
\node[fill=quotientorange!30] (s) at (0,0) {$s$};
\draw[targetgreen,very thick] (t1)--(s) (t2)--(s) (t3)--(s) (t4)--(s);
\draw[edgegray,thick,dashed] (t1)--(t2) (t2)--(t3) (t3)--(t4) (t4)--(t1);
\node[draw=none,font=\scriptsize,targetgreen,right=1.1cm of t2,align=left,yshift=-0.4cm]
  {star via $s$: $4\times 2 = 8$};
\node[draw=none,font=\scriptsize,right=1.1cm of t3,align=left,yshift=0.4cm]
  {tree on terminals only:\\$3$ sides $\times\,3 = 9$};
\end{tikzpicture}
\caption{Why an optional vertex helps. Terminals $t_1,\dots,t_4$ (blue) must be
connected; $s$ (orange) is a non-terminal that may be recruited. Solid green
spokes cost $2$ each, dashed perimeter edges cost $3$ each. The Steiner star
through $s$ costs $8$; the cheapest tree using only terminal edges needs three
perimeter sides and costs $9$. Choosing \emph{whether} to include $s$ is the
combinatorial decision that makes the general problem hard.}
\label{ch37:fig-steiner}
\end{figure}

\section{A Map of the Reductions}
\label{ch37:map}

\begin{figure}[H]
\centering
\begin{tikzpicture}[
  box/.style={rectangle,rounded corners,draw,thick,fill=blue!5,
    font=\footnotesize,minimum height=7mm,inner sep=3pt,align=center},
  anchorbox/.style={rectangle,rounded corners,draw,very thick,fill=yellow!18,
    font=\footnotesize\bfseries,minimum height=8mm,inner sep=3pt},
  ar/.style={-{Stealth[length=2.2mm]},thick,structurepurple}]
\node[anchorbox] (sat) at (0,0) {3-SAT};
\node[box] (nae)  at (-4.6,0) {NAE-3SAT};
\node[box] (mc)   at (-5.6,-1.6) {Max-Cut};
\node[box] (is)   at (-2.6,-1.6) {Independent\\Set};
\node[box] (col)  at (0.4,-1.6) {3-Coloring};
\node[box] (ham)  at (3.3,-1.6) {Ham.\\Cycle};
\node[box] (long) at (5.7,-1.6) {Longest\\Path};
\node[box] (vc)   at (-4.0,-3.2) {Vertex\\Cover};
\node[box] (cl)   at (-1.4,-3.2) {Clique};
\node[box] (col4) at (0.8,-3.2) {$c$-Coloring};
\node[box] (tsp)  at (3.3,-3.2) {TSP};
\node[box] (st)   at (-4.0,-4.6) {Steiner\\Tree};
\draw[ar] (sat)--(is);
\draw[ar] (sat)--(col);
\draw[ar] (sat)--(ham);
\draw[ar] (sat)--(nae);
\draw[ar] (nae)--(mc);
\draw[ar] (is)--(vc);
\draw[ar] (is)--(cl);
\draw[ar] (col)--(col4);
\draw[ar] (ham)--(tsp);
\draw[ar] (ham)--(long);
\draw[ar] (vc)--(st);
\end{tikzpicture}
\caption{The graph wing of Karp's web. Every arrow is a polynomial reduction
proved or sketched in this chapter; hardness flows along the arrows out of the
3-SAT anchor. A polynomial algorithm for any single box would imply
$\mathsf{P}=\mathsf{NP}$.}
\label{ch37:fig-map}
\end{figure}

\section{The Boundary Between Easy and Hard}
\label{ch37:boundary}

The most instructive part of this subject is
\emph{where} the line between easy and hard falls. Tiny changes in a problem statement -- $2$
versus $3$, ``edge'' versus ``vertex,'' shortest versus longest -- flip
tractability. We look at five boundary cases and the reason for each.

\subsection{2-Coloring vs 3-Coloring}

\textsc{2-Coloring} is polynomial: a graph is 2-colorable iff it is
\emph{bipartite} iff it has no odd cycle, all checkable by a single BFS/DFS that
2-colors greedily and reports failure on a conflicting edge
(Figure~\ref{ch37:fig-2col}).

\begin{algorithm}[H]
\caption{\textsc{Two-Color}$(\Gph)$ -- bipartiteness by BFS}
\label{ch37:alg-2col}
\DontPrintSemicolon
\KwIn{Undirected graph $\Gph=(V,E)$.}
\KwOut{A proper 2-coloring, or \textsf{not 2-colorable}.}
$\mathrm{color}[v]\gets\textsf{none}$ for all $v$\;
\ForEach{vertex $s$ with $\mathrm{color}[s]=\textsf{none}$}{
  $\mathrm{color}[s]\gets 0$;\quad enqueue $s$\;
  \While{queue nonempty}{
    $u\gets$ dequeue\;
    \ForEach{neighbor $w$ of $u$}{
      \uIf{$\mathrm{color}[w]=\textsf{none}$}{
        $\mathrm{color}[w]\gets 1-\mathrm{color}[u]$;\quad enqueue $w$\;
      }
      \ElseIf{$\mathrm{color}[w]=\mathrm{color}[u]$}{
        \Return \textsf{not 2-colorable}\;
      }
    }
  }
}
\Return $\mathrm{color}$\;
\end{algorithm}

\begin{whybox}
The algorithm is forced at every step: once $s$ is color $0$, all neighbors
\emph{must} be $1$, theirs $0$, and so on -- the colors propagate with zero
freedom. A conflict (an edge whose endpoints land on the same color) is reachable
iff the BFS tree plus that edge forms an odd cycle, and odd cycles are precisely
the obstruction to bipartiteness. With \emph{three} colors the propagation is no
longer forced: at a vertex with one colored neighbor you have two free choices,
and those branching choices are exactly the combinatorial explosion that 3-SAT
injects through the gadget of \S\ref{ch37:coloring}. Two colors leave no slack;
three colors leave enough slack to encode arbitrary computation.
\end{whybox}

\begin{figure}[H]
\centering
\begin{tikzpicture}[every node/.style={circle,draw,thick,inner sep=1.5pt,
  font=\scriptsize,minimum size=15pt},scale=1.0]
\node[fill=green!18] (p1) at (0,1.4) {$0$};
\node[fill=red!18]   (p2) at (1.3,0.6) {$1$};
\node[fill=green!18] (p3) at (0.8,-0.9) {$0$};
\node[fill=red!18]   (p4) at (-0.8,-0.9) {$1$};
\node[fill=green!18] (p5) at (-1.3,0.6) {$0$};
\draw[edgegray,thick] (p1)--(p2)--(p3)--(p4)--(p5);
\draw[trapred,thick] (p5)--(p1);
\node[draw=none,trapred,right=1.6cm of p2,font=\scriptsize,align=left]
  {closing edge $p_5p_1$:\\both color $0$ -- conflict};
\end{tikzpicture}
\caption{\textsc{Two-Color} on the 5-cycle. BFS from $p_1$ forces the colors
$0,1,0,1,0$ around the path $p_1p_2p_3p_4p_5$; the closing edge $p_5p_1$ then
joins two $0$'s, so Algorithm~\ref{ch37:alg-2col} reports \textsf{not
2-colorable}. The obstruction is exactly the odd cycle.}
\label{ch37:fig-2col}
\end{figure}

\subsection{2-SAT vs 3-SAT}

\textsc{2-SAT} (two literals per clause) is polynomial; \textsc{3-SAT} is
$\mathsf{NP}$-complete. The difference is structural: a 2-clause $(a\vee b)$ is
equivalent to two implications $\bar a\Rightarrow b$ and $\bar b\Rightarrow a$.
Build the \emph{implication graph} on the $2n$ literals; then $\varphi$ is
satisfiable iff no variable $x_i$ has $x_i$ and $\bar x_i$ in the same strongly
connected component. That is a single SCC computation (Chapter 5), $\Oh(n+m)$.

\begin{intuitionbox}
A 2-clause forces a chain reaction: setting one literal cascades
deterministically through implications, like the 2-coloring propagation. A
3-clause $(a\vee b\vee c)$ gives only the weak implication $(\bar a\wedge\bar
b)\Rightarrow c$ -- you must already know \emph{two} values before it constrains
the third. That ``need two to pin one'' is genuine disjunctive choice, and choice
is what makes search exponential. The recurring pattern across all five
boundaries: hardness appears exactly when local constraints stop forcing unique
propagation and start offering free binary choices.
\end{intuitionbox}

\subsection{Shortest vs Longest Path}

Covered in \S\ref{ch37:longest}: shortest paths have optimal substructure
(prefixes of optimal solutions are optimal) because nonnegative cycles never
help; longest \emph{simple} paths lose substructure to the simplicity
constraint. On DAGs, where simplicity is automatic, longest path is polynomial
again -- the boundary is the presence of cycles.

\subsection{Min-Cut vs Max-Cut}

Min-Cut is polynomial (max-flow duality, Chapter 12; Karger's randomized global
cut, Chapter 11). Max-Cut is $\mathsf{NP}$-hard (\S\ref{ch37:maxcut-steiner}).
Minimizing a cut is a flow problem with LP duality and integral optima;
maximizing the cut is equivalent to a $\pm1$ quadratic optimization with no such
duality -- the feasible region's geometry is convex for min and combinatorial for
max.

\subsection{Matching vs 3-Dimensional Matching}

Maximum matching in graphs is polynomial -- Edmonds' blossom algorithm (Chapter
17). \textsc{3-Dimensional-Matching} (perfectly pair up triples from three
disjoint sets) is $\mathsf{NP}$-complete (Karp). Two-way matching has the
augmenting-path / LP-integrality structure; the three-way version does not, and
it is itself a primary source for reductions to Steiner tree, set cover, and
scheduling.

\begin{table}[H]
\centering\small
\begin{tabular}{lll}
\toprule
\textbf{Easy (in $\mathsf{P}$)} & \textbf{Hard ($\mathsf{NP}$-complete)} & \textbf{What flips it}\\
\midrule
2-Coloring (bipartite) & 3-Coloring & forced vs free propagation\\
2-SAT & 3-SAT & implication vs disjunction\\
Shortest path & Longest simple path & cycles + simplicity\\
Eulerian cycle & Hamiltonian cycle & edges vs vertices\\
Min-cut & Max-cut & flow duality vs $\pm1$ quadratic\\
Matching ($2$D) & $3$D matching & augmenting paths vs none\\
MST ($R=V$) / SP ($\abs R=2$) & Steiner tree & optional-vertex choice\\
\bottomrule
\end{tabular}
\caption{The easy/hard frontier. Each ``hard'' problem is one small modification
away from a polynomial sibling; the right column names the structural feature
that disappears.}
\label{ch37:tab-boundary}
\end{table}

\section{What Hardness Means for Algorithm Design}
\label{ch37:design}

A proof that your problem is $\mathsf{NP}$-hard is useful
guidance: it tells you to stop looking for an exact polynomial
algorithm and pick one of the legitimate escape routes, each developed in its own
later chapter.

\begin{itemize}
  \item \textbf{Approximation} (Chapter 38). Settle for provably near-optimal.
        Vertex-Cover has a $2$-approximation (take both endpoints of a maximal
        matching); metric TSP has Christofides' $3/2$; Max-Cut has
        Goemans-Williamson's $0.878$. Some problems resist: general TSP and
        \textsc{Clique} have no good approximation unless
        $\mathsf{P}=\mathsf{NP}$ (the latter by the PCP theorem).
  \item \textbf{Parameterization} (Chapters 24--25). Confine the exponential
        blow-up to a parameter $k$ that is small in practice. \textsc{Vertex-Cover}
        is fixed-parameter tractable in $\Oh(1.27^k + kn)$ -- a
        bounded-search-tree algorithm -- so it is solvable exactly when the cover
        is small even if the graph is huge. \textsc{Clique} and
        \textsc{Independent-Set}, by contrast, are $\mathsf{W}[1]$-hard: no such
        $f(k)\cdot\poly(n)$ algorithm is expected.
  \item \textbf{Structural restriction.} Many hard problems collapse to
        polynomial on bounded-treewidth graphs (Chapter 24), planar graphs
        (Chapter 26; Baker's technique gives PTASs), or interval/chordal/perfect
        graphs (Chapter 30, where coloring and clique become polynomial). If your
        inputs have structure, exploit it.
  \item \textbf{Exact exponential algorithms.} When $n$ is modest, an $\Oh(2^n)$
        Held-Karp DP solves TSP, and branch-and-bound with good pruning routinely
        handles thousands of cities. Hardness is asymptotic; small instances are
        fair game.
  \item \textbf{Heuristics and SAT/ILP solvers.} Encode back into SAT or integer
        programming and run an industrial solver. Modern CDCL SAT solvers and
        MIP engines solve very large structured instances despite worst-case
        hardness.
\end{itemize}

\begin{aestheticbox}
The main lesson of this chapter is how closely these problems are linked. Nine graph problems that
arose in nine different decades and applications -- routing, scheduling, register
allocation, network design -- turned out to be computationally equivalent,
linked together by polynomial-time reductions. Proving one hard proves them all hard; finding one
easy would make them all easy. That web is why ``$\mathsf{NP}$-complete''
carries so much weight: it is a structural
claim about thousands of problems at once. When you recognize your problem as a
member, you inherit fifty years of accumulated knowledge about how to cope.
\end{aestheticbox}

\begin{connectionbox}
This chapter is the graph-flavored half of the $\mathsf{NP}$-completeness chapter
in the companion algorithms book. That book proves Cook-Levin from the
Turing-machine tableau and develops SAT, subset-sum, and partition; here we took
3-SAT as a black box and grew the graph wing -- Independent Set, Coloring,
Hamiltonicity, TSP, Longest Path, Max-Cut, Steiner Tree -- by gadget reductions.
The escape routes (approximation, FPT, structural restriction) are the subject of
Chapters 24, 25, 26, 30, and 38. Max-Cut in particular threads forward into the
spectral and Laplacian chapters (31--32), where its relaxation is studied in terms of
eigenvalues and electrical flows.
\end{connectionbox}

\section{Historical Context and Controversies}

\begin{historybox}
The mathematical realization that certain computational problems might be intrinsically intractable began with Kurt G\"odel. In a $1956$ letter to John von Neumann, G\"odel asked whether finding mathematical proofs of length $n$ could be decided mechanically in $\Oh(n)$ or $\Oh(n^2)$ steps, noting that if so, the discovery of mathematical truth could be completely automated. Throughout the $1950$s and $1960$s, researchers in operations research repeatedly got stuck: while linear programming, spanning trees, and shortest paths yielded to polynomial algorithms, seemingly modest variations -- such as requiring integer solutions, finding a traveling salesman tour, or seeking a maximum clique -- stumped every algorithmic attempt.

In $1971$, Stephen Cook formalized the class $\mathbf{NP}$ and proved that Boolean Satisfiability (\textsc{SAT}, deciding whether a propositional formula can be made true) is $\mathbf{NP}$-complete, showing that every problem whose solutions can be verified in polynomial time reduces to \textsc{SAT}. Working independently in the Soviet Union, Leonid Levin published similar results in $1973$ on universal search problems. The full impact of Cook's insight was demonstrated by Richard Karp in $1972$, who proved that $21$ central combinatorial problems -- including \textsc{Vertex Cover}, \textsc{Clique}, \textsc{Hamiltonian Cycle}, \textsc{Coloring}, and \textsc{Max-Cut} -- are all $\mathbf{NP}$-complete, showing that computational hardness is widespread in discrete mathematics rather than an oddity of logic.

Before Cook and Karp, researchers spent careers searching for polynomial-time algorithms for these problems, attributing failure to personal shortcoming rather than to a mathematical obstacle. A persistent trap is confusing the hardness of \emph{search} with the ease of \emph{verification}: checking a proposed certificate takes polynomial time, but finding it through exhaustive search requires exploring exponentially many configurations. Proving $\mathbf{NP}$-completeness redirected theoretical energy toward approximation, fixed-parameter tractability, and heuristics. The main open question in theoretical computer science remains the $\mathbf{P}$ versus $\mathbf{NP}$ problem itself: proving mathematically that search is strictly harder than verification.
\end{historybox}

\subsection{What We Can Now Do and Its Limitations}
\begin{connectionbox}
\textbf{Capabilities Unlocked}:
\begin{itemize}
  \item We have a general method for proving that a graph problem is $\mathbf{NP}$-hard. A polynomial reduction from a known $\mathbf{NP}$-complete problem proves that finding an exact polynomial algorithm is as hard as solving every problem in $\mathbf{NP}$.
  \item Proving $\mathbf{NP}$-hardness tells algorithm designers to stop searching for exact polynomial algorithms and turn to practical alternatives: approximation algorithms, parameterized complexity ($\mathbf{FPT}$), and structural exploitation (e.g., planarity, bounded treewidth).
\end{itemize}

\textbf{Limitations}:
\begin{itemize}
  \item \emph{Worst-case pessimism}: $\mathbf{NP}$-completeness is strictly a worst-case classification. A problem can be $\mathbf{NP}$-hard even though $99\%$ of real-world instances encountered in practice can be solved in seconds by modern SAT solvers or integer programming packages.
  \item \emph{Reduction discipline}: Beginners frequently execute reductions in reverse. Reducing a known polynomial problem (such as Max-Flow) to an unknown problem proves nothing about the unknown problem's hardness; hardness flows strictly \emph{from} the known hard problem \emph{to} the target problem.
  \item \emph{Destruction of approximation ratios}: An exact polynomial reduction between two decision problems generally preserves only exact solvability, not approximate structure. For example, while Hamiltonian Cycle reduces to General TSP, Metric TSP admits a $3/2$-approximation, whereas General TSP cannot be approximated within any constant factor unless $\mathbf{P} = \mathbf{NP}$.
\end{itemize}

\textbf{Where We Go Next}:
\begin{itemize}
  \item In Chapter~38 (Approximation Algorithms on Graphs), we trade exactness for provable worst-case guarantees in polynomial time, developing LP rounding, greedy set cover, Christofides' $3/2$-approximation for metric TSP, and Goemans--Williamson's $0.878$-approximation for Max-Cut via semidefinite programming.
  \item In Chapter~39 (Dynamic Graph Algorithms), we study how to maintain connectivity, spanning forests, and shortest-path structures under online edge insertions and deletions without recomputing from scratch.
\end{itemize}
\end{connectionbox}

\section{Exercises}
\label{ch37:exercises}

\begin{exercise}[Direction discipline]
\label{ch37:ex-direction}
A classmate ``proves'' \textsc{Clique} is easy by reducing it to
\textsc{Max-Flow}: ``find a clique by computing a max-flow, which is
polynomial.'' Explain in one paragraph exactly which arrow points the wrong way,
and why a reduction \emph{from} Clique \emph{to} a polynomial problem proves
nothing about Clique's difficulty.
\end{exercise}

\begin{exercise}[Complementarity by hand]
\label{ch37:ex-complement}
Take the 5-cycle $C_5$. (a) Find its maximum independent set and its minimum
vertex cover, and verify $\abs{\text{IS}}+\abs{\text{VC}}=5$. (b) Draw
$\overline{C_5}$ and find its maximum clique; confirm it equals the maximum
independent set of $C_5$ via Lemma~\ref{ch37:lem-triangle}.
\end{exercise}

\begin{exercise}[Run the IS gadget]
\label{ch37:ex-isgadget}
Apply Algorithm~\ref{ch37:alg-is} to
$\varphi=(x_1\vee x_2\vee \bar x_3)\wedge(\bar x_1\vee x_2\vee x_3)\wedge
(x_1\vee \bar x_2\vee x_3)$. Draw the graph (three triangles plus conflict edges),
find an independent set of size $3$, and read off the satisfying assignment.
\end{exercise}

\begin{exercise}[Why exactly one per triangle]
\label{ch37:ex-trianglecount}
In the proof of Theorem~\ref{ch37:thm-is}, we asserted that an independent set of
size $m$ must take exactly one vertex from each clause triangle. Prove it
carefully: show that ``at most one per triangle'' plus ``$m$ triangles'' plus
``$\abs S=m$'' forces ``exactly one per triangle,'' and explain where the
argument would break if some clause had only two literals.
\end{exercise}

\begin{exercise}[Metric vs general TSP]
\label{ch37:ex-tsp}
(a) Verify that the $1$/$2$ weighting in Theorem~\ref{ch37:thm-tsp} satisfies the
triangle inequality. (b) Show that if general (non-metric) TSP had a polynomial
$\rho$-approximation for any constant $\rho\ge 1$, then
$\mathsf{P}=\mathsf{NP}$, using the heavy-non-edge reduction discussed after
Theorem~\ref{ch37:thm-tsp}. (c) Why does the same argument fail to rule out the Christofides
$3/2$-approximation for metric TSP?
\end{exercise}

\begin{exercise}[Longest path on a DAG]
\label{ch37:ex-dag}
Give an $\Oh(V+E)$ algorithm for the longest path in a DAG. Then exhibit a small
non-DAG (say, a graph on $5$ vertices) where your DP returns a wrong answer, and
pinpoint the step that assumed acyclicity.
\end{exercise}

\begin{exercise}[2-SAT implication graph]
\label{ch37:ex-2sat}
Solve $(x_1\vee x_2)\wedge(\bar x_1\vee x_3)\wedge(\bar x_2\vee\bar x_3)\wedge
(x_2\vee x_3)$ by building the implication graph on the six literals, computing
its strongly connected components, and checking that no $x_i$ shares a component
with $\bar x_i$. Read off a satisfying assignment from the condensation's reverse
topological order.
\end{exercise}

\begin{exercise}[Coloring lifts 3 to $c$]
\label{ch37:ex-clift}
Prove that \textsc{$c$-Coloring} is $\mathsf{NP}$-complete for every fixed $c\ge
3$ by reducing \textsc{3-Coloring}: add $c-3$ new vertices forming a clique, each
joined to every original vertex, and argue that the augmented graph is
$c$-colorable iff the original is $3$-colorable.
\end{exercise}

\begin{exercise}[Break a gadget]
\label{ch37:ex-break}
In the 3-coloring reduction (\S\ref{ch37:coloring}), suppose you \emph{forgot} to
make each clause-output vertex adjacent to the base color $B$ (you only tied it to
$F$). Construct a small unsatisfiable formula whose graph is nevertheless
3-colorable under the broken gadget, demonstrating that the reduction is no longer
correct.
\end{exercise}

\begin{exercise}[FPT is not free]
\label{ch37:ex-fpt}
\textsc{Vertex-Cover} parameterized by solution size $k$ has an $\Oh(2^k(n+m))$
bounded-search-tree algorithm: pick any uncovered edge $uv$ and branch on putting
$u$ or $v$ in the cover. (a) Write this as pseudocode and prove correctness and
the $\Oh(2^k(n+m))$ bound. (b) Explain why the complementary trick
($\text{IS}=n-\text{VC}$) does \emph{not} yield an $f(k)\cdot\poly(n)$ algorithm
for \textsc{Independent-Set} parameterized by the size of the independent set.
\end{exercise}

\begin{exercise}[Max-Cut on a triangle and a square]
\label{ch37:ex-maxcut}
(a) Compute the max-cut of $K_3$ and of $K_4$ by hand. (b) For the $4$-cycle
$C_4$, exhibit a cut of all $4$ edges and argue it is optimal. (c) Sketch why
greedy local search (move any vertex that increases the cut) terminates and
reaches a cut of size $\ge \abs E/2$, a $2$-approximation -- and give a small
graph where it gets stuck below optimal.
\end{exercise}

\begin{exercise}[Reduction transitivity in practice]
\label{ch37:ex-transitive}
Compose the chapter's reductions to give a \emph{single direct} polynomial map
from \textsc{3-SAT} to \textsc{TSP}, and state the resulting bound $L$ and number
of cities in terms of the number of variables $n$ and clauses $m$. Confirm the
composed map is still polynomial.
\end{exercise}
